\chapter{Introduction}
\label{ch:introduction}

Professional football is a global industry, and the transfer market is among its most visible parts.\footnote{Generative AI tools were used extensively in this thesis, for code, for first drafts of the text and for review. Appendix~\ref{app:ai} states how.} In 2019, 18,080 players moved between clubs in different countries and clubs paid USD~7.35~billion to recruit them \cite{franceschi2024determinants}. Decisions of this size are still made with a good deal of judgement. Clubs overpay for players who underperform and overlook players whose value conventional scouting does not capture. Data analytics has therefore become a tool to support more objective and replicable estimates of player value \cite{muller2017beyond}.

An active strand of football analytics profiles players by their performance attributes, to support scouting, squad construction and tactical planning. Clustering is particularly useful here because it groups players into profiles without relying on positional labels, which capture the role of a modern player only roughly. \citeA{durso2022robust} propose a robust fuzzy $C$-medoids model for mixed-type data that handles counts, success rates and categorical positions together, assigns each player a graded degree of membership in each profile, and puts outliers into a separate noise cluster. \citeA{akhanli2023clustering} develop mixed-type dissimilarities and validity indices for the same task. Together these contributions make fuzzy clustering of mixed-type performance data a rigorous way of profiling players.

Almost all of this work is cross-sectional, however. Players are clustered in a single season, so the result is a snapshot. Player development, by contrast, is dynamic: roles shift with age, club changes alter tactical contexts, and similar players follow different paths. The longitudinal literature is scarce. The closest study, by \citeA{pacifico2025fuzzy}, analyses 27 representative players over twelve seasons and does not relate profile changes to market values. \citeA{durso2022robust} themselves point to the temporal dimension of the playing variables and to financial variables as natural extensions. Whether the way in which a player's profile \emph{changes} from one season to the next is associated with a change in his market value, once his current performance is taken into account, has not been studied. This is the question of the thesis.

\section{Research questions and approach}

The main research question is:
\begin{quote}
\emph{Do longitudinal transitions in fuzzy cluster membership, derived from a robust mixed-data clustering model applied to Primeira Liga player performance statistics across consecutive seasons, have incremental explanatory power for changes in transfer market valuations beyond what individual performance metrics and basic player and club characteristics capture?}
\end{quote}
Two supporting questions are descriptive and methodological. The first asks which archetypes the clustering identifies, how stable they are across seasons, and how players move between them. The second asks how sensitive the answer to the main question is to the choices that have to be made when the transitions are constructed. The objective is explanatory and associational. The design is observational, no causal identification strategy is introduced, and all findings are reported as associations.

The thesis applies the model of \citeA{durso2022robust} to eight seasons (2017/18 to 2024/25) of the Portuguese Primeira Liga. Each season is clustered separately, the archetypes are aligned across seasons with the Hungarian algorithm, and the change in each player's membership vector between consecutive seasons defines a transition. The change in the logarithm of the player's Transfermarkt market value over the same period is then related to the transition in a regression with controls for the established determinants of value. Because the memberships are estimated, inference is based on a two-stage bootstrap that repeats the clustering and the regression \cite{pagan1984econometric}. The Primeira Liga is a deliberate choice. The league supplies many players to the larger leagues, yet it has received little attention in quantitative research on player valuation, which \citeA{franceschi2024determinants} identify as a gap.

\section{Findings in brief}

The clustering yields four interpretable archetypes of outfield players (defenders, central midfielders, wide attackers and strikers) and a small noise cluster of high-volume attackers. From one season to the next, 75\% of the players remain in their archetype. The estimation of the attribute weights, which is a distinguishing feature of the model, failed on these data and was replaced by weights fixed in proportion to the number of variables in each group. The archetypes had to be aligned across seasons by their profiles instead of their medoids. In the main configuration, which was fixed before the valuation stage was estimated, membership transitions are weakly associated with changes in market value beyond the two baselines, with small gains in explanatory power. The association is concentrated in the shift from the defender to the wide-attacker profile. The association is fragile: it is absent for some settings of the clustering, it varies strongly between random draws of the clustering, and no result survives a correction for the number of specifications examined. It also largely disappears when changes in the position labels, which enter the clustering but not the baselines, are controlled. The thesis therefore finds no robust incremental explanatory power. The alignment method and the data were corrected, and the position check was added, after earlier valuation results had been seen. Sections~\ref{sec:rob-alignment} and~\ref{sec:results-position} report this.

\section{Relevance}

Scientifically, the thesis adds a longitudinal dimension to the fuzzy clustering of mixed-type football data and links it to market valuation, so it combines two literatures that have developed separately. It documents a failure of weight estimation in the model of \citeA{durso2022robust} and a workable alternative, and it shows how sensitive an application of fuzzy clustering can be to its tuning. For practice, the question is whether the trajectory of a player's profile carries information about his value that scouts and sporting directors can use. The answer given here is cautious. The archetypes are a useful description, but their changes are not a reliable basis for valuation.

\section{Structure}

Chapter~\ref{ch:literature} reviews the literature on player valuation, fuzzy clustering and the clustering of football players, and states the research questions. Chapter~\ref{ch:theory} develops the theoretical argument and the hypothesis. Chapter~\ref{ch:data} describes the data and Chapter~\ref{ch:methodology} the methods. Chapter~\ref{ch:results} reports the results for the main configuration, and Chapter~\ref{ch:robustness} their sensitivity. Chapter~\ref{ch:discussion} discusses the findings, implications and limitations, and Chapter~\ref{ch:conclusion} concludes.
